% =========================================================
% Order Axioms
% =========================================================

\section{Order Axioms}
\label{sec:order-axioms-q}

The rational numbers are not merely a field. They also carry an order
relation compatible with the algebraic operations, so that $\mathbb{Q}$
becomes an ordered field.

\begin{remark}
For many arguments it is cleaner to encode order using the set of
positive elements rather than to take $<$ as primitive.
\end{remark}

\begin{definition}[Positive cone]
\label{def:positive-cone}
Let $F$ be a field. A subset $P\subset F$ is called a \emph{positive
cone} if it satisfies:
\begin{enumerate}
  \item if $a,b\in P$, then $a+b\in P$;
  \item if $a,b\in P$, then $ab\in P$;
  \item for every $a\in F$, exactly one of the following holds:
  \[
  a\in P,
  \qquad
  a=0,
  \qquad
  -a\in P.
  \]
\end{enumerate}
Define the order relation by
\[
a<b \iff b-a\in P.
\]
\end{definition}

\begin{remark}
This viewpoint makes the interaction between algebra and order
transparent. Inequalities can be translated into statements about
positive elements and then handled algebraically.
\end{remark}

\begin{definition}[Ordered integral domain]
\label{def:ordered-integral-domain}
An \emph{ordered integral domain} is an integral domain
\[
(F,+,\cdot,0,1)
\]
together with a relation $<$ on $F$ such that:
\begin{enumerate}
  \item for all $a,b\in F$, exactly one of the relations
  \[
  a<b,\qquad a=b,\qquad b<a
  \]
  holds;
  \item if $a<b$, then for every $c\in F$,
  \[
  a+c<b+c;
  \]
  \item if $a<b$ and $0<c$, then
  \[
  ac<bc.
  \]
\end{enumerate}
\end{definition}

\begin{definition}[Ordered field]
\label{def:ordered-field}
An \emph{ordered field} is a field
\[
(F,+,\cdot,0,1)
\]
together with a relation $<$ on $F$ making $F$ into an ordered integral
domain.
\end{definition}

\begin{definition}[Ordered subfield]
\label{def:ordered-subfield}
Let $E$ and $F$ be ordered fields with $E\subseteq F$. We say that $E$ is an
\emph{ordered subfield} of $F$ if $E$ is a subfield of $F$ and the order on
$E$ is the restriction of the order on $F$.
\end{definition}

\begin{definition}[Densely ordered system]
\label{def:densely-ordered-system}
Let $S$ be a set with a relation $<$. We say that $(S,<)$ is
\emph{densely ordered} if:
\begin{enumerate}
  \item $(S,<)$ is a simply ordered system; and
  \item for every $x,y\in S$ with $x<y$, there exists $z\in S$ such that
  \[
  x<z<y.
  \]
\end{enumerate}
\end{definition}

\begin{theorem}
\label{thm:q-is-ordered-field}
With its usual order, the system
\[
(\mathbb{Q},+,\cdot,0,1,<)
\]
is an ordered field.
\end{theorem}

\begin{remark*}[Proof]
\hyperref[prf:q-is-ordered-field]{Go to proof.}
\end{remark*}

\begin{remark}
This theorem records the compatibility of the rational order with the field
operations. It is this compatibility that makes inequalities in $\mathbb{Q}$
behave as expected under addition, subtraction, multiplication, and division
by positive elements.
\end{remark}

\begin{theorem}
\label{thm:squares-nonnegative}
For every $a\in F$,
\[
a^2\ge 0.
\]
\end{theorem}

\begin{remark*}[Proof]
\hyperref[prf:squares-nonnegative]{Go to proof.}
\end{remark*}

\begin{theorem}
\label{thm:inverse-positive}
If $a>0$, then $a^{-1}>0$.
\end{theorem}

\begin{remark*}[Proof]
\hyperref[prf:inverse-positive]{Go to proof.}
\end{remark*}

\begin{theorem}
\label{thm:multiply-positive}
If $a<b$ and $c>0$, then
\[
ac<bc.
\]
\end{theorem}

\begin{remark*}[Proof]
\hyperref[prf:multiply-positive]{Go to proof.}
\end{remark*}

\begin{theorem}
\label{thm:multiply-negative}
If $a<b$ and $c<0$, then
\[
ac>bc.
\]
\end{theorem}

\begin{remark*}[Proof]
\hyperref[prf:multiply-negative]{Go to proof.}
\end{remark*}

\begin{theorem}
\label{thm:one-positive}
In an ordered field,
\[
1>0.
\]
\end{theorem}

\begin{remark*}[Proof]
\hyperref[prf:one-positive]{Go to proof.}
\end{remark*}

\begin{corollary}
\label{cor:n-positive}
For every natural number $n$,
\[
n>0.
\]
\end{corollary}

\begin{remark*}[Proof]
\hyperref[prf:n-positive]{Go to proof.}
\end{remark*}

\begin{theorem}
\label{thm:positive-fraction-iff-positive-product}
For any $a,b\in F$ with $b\neq 0$,
\[
0<\frac{a}{b}
\quad\Longleftrightarrow\quad
0<a\cdot b.
\]
\end{theorem}

\begin{remark*}[Proof]
\hyperref[prf:positive-fraction-iff-positive-product]{Go to proof.}
\end{remark*}

\begin{theorem}
\label{thm:fraction-order-comparison-general}
For any $a,b,c,d\in F$ with $b\neq 0$ and $d\neq 0$,
\[
\frac{a}{b}<\frac{c}{d}
\quad\Longleftrightarrow\quad
(a\cdot d)\cdot(b\cdot d)<(b\cdot c)\cdot(b\cdot d).
\]
\end{theorem}

\begin{remark*}[Proof]
\hyperref[prf:fraction-order-comparison-general]{Go to proof.}
\end{remark*}

\begin{theorem}
\label{thm:fraction-order-comparison-positive-denominator-product}
For any $a,b,c,d\in F$, if $b\neq 0$, $d\neq 0$, and $0<b\cdot d$, then
\[
\frac{a}{b}<\frac{c}{d}
\quad\Longleftrightarrow\quad
a\cdot d<b\cdot c.
\]
\end{theorem}

\begin{remark*}[Proof]
\hyperref[prf:fraction-order-comparison-positive-denominator-product]{Go to proof.}
\end{remark*}

\begin{theorem}
\label{thm:reciprocal-order-for-positive-elements}
If $0<d<b$, then
\[
0<\frac{1}{b}<\frac{1}{d}.
\]
Conversely, if
\[
0<\frac{1}{b}<\frac{1}{d},
\]
then $0<d<b$.
\end{theorem}

\begin{remark*}[Proof]
\hyperref[prf:reciprocal-order-for-positive-elements]{Go to proof.}
\end{remark*}

\begin{theorem}
\label{thm:ordered-fields-are-densely-ordered}
Every ordered field is densely ordered by its order relation.
\end{theorem}

\begin{remark*}[Proof]
\hyperref[prf:ordered-fields-are-densely-ordered]{Go to proof.}
\end{remark*}

\begin{remark}[Transition]
\label{rem:arch-v-complete}
The positivity of $1$ anchors the order structure and leads naturally to
the Archimedean property. That property explains why natural numbers
eventually outrun every fixed rational and why reciprocals like $1/n$
become arbitrarily small.
\end{remark}
